\documentclass[14pt,a4paper]{extreport}
% preamble.tex — спільна преамбула для всіх лабораторних робіт. Підключається одразу після \documentclass у кожному lab*/main.tex.
%
% ЗБІРКА: LuaLaTeX (не pdfLaTeX). Причина: пакет listings під pdfLaTeX не читає UTF-8 і зупиняє збірку на кириличних коментарях у лістингах програмного коду. LuaLaTeX обробляє UTF-8 нативно.
%
% ШРИФТ: Liberation Serif — метрично сумісний з Times New Roman і містить повний набір кириличних гліфів (TeX Gyre Termes кирилиці не має, через що український текст не відображався взагалі).

\usepackage{fontspec}
\setmainfont{Liberation Serif}
\setmonofont{DejaVu Sans Mono}[Scale=0.82]

\usepackage{polyglossia}
\setmainlanguage{ukrainian}
\setotherlanguage{english}

\usepackage{microtype}

% Поля, інтервал, відступи й заголовки нижче взяті з шаблону НУЛП, за яким написана бакалаврська робота (../bachelor-thesis/thesis/main.tex). Той самий формат приймає В.М. Шиманський і для магістерських, тож усі звіти тримаємо в ньому, а не в тому, що переказували на консультації.
\usepackage[left=25mm,right=15mm,top=15mm,bottom=23mm,footskip=1.5mm]{geometry}
\usepackage{setspace}
% Саме \setstretch{1.5}, а не \onehalfspacing: останній дає близько 1.25.
\setstretch{1.5}

\usepackage{graphicx}
\usepackage{float}
\usepackage{booktabs}
\usepackage{array}
\usepackage{tabularx}
\usepackage{longtable}
\usepackage{enumitem}
% labelindent вирівнює маркер списку з абзацним відступом — так у шаблоні НУЛП.
\setlist{nosep}
\setlist[itemize]{label=\textbullet,leftmargin=*,labelindent=1.27cm}
\setlist[enumerate]{leftmargin=*,labelindent=1.27cm}
\usepackage{indentfirst}
\usepackage{amsmath}

\usepackage{hyperref}
\hypersetup{colorlinks=true,linkcolor=black,urlcolor=blue,citecolor=black}

% xurl дозволяє \url розривати адресу в будь-якому місці, а не лише після скісної риски. Без нього довгий сегмент шляху (напр. «scopus---podrobnaya- instruktsiya-dlya-uchenogo») виходить за межі сторінки, бо дефіс для hyperref не є точкою розриву. Вантажиться після hyperref — перевизначає його \UrlBreaks.
\usepackage{xurl}

% ---------------------------- Список використаної літератури ----------------------------
% biblatex + biber зі стилем IEEE: посилання «[1]» у тексті, записи — за порядком першої згадки. Обрано IEEE, а не gost-numeric (ДСТУ 8302:2015), — так вирішив автор звітів після порівняння обох варіантів.
% Кожен звіт із джерелами сам підключає свій references.bib через \addbibresource і друкує список \printbibliography[heading=none] під власним \section. Звіт без джерел від завантаженого biblatex нічого не втрачає.
% Оточення списку — стандартне для стилю: мітка від лівого поля, текст опису вирівняний після неї. Не підганяємо його під абзацний відступ нумерованих переліків тексту (labelindent у \setlist вище) — список літератури є окремим розділом зі своїм стандартним оформленням.
\usepackage[backend=biber,style=ieee]{biblatex}
% biblatex за замовчуванням додає розтяжний проміжок після «/ . : ? =» в адресах, і з вимкненими переносами рядок вирівнюється розрідженням самої адреси («https : / / www . scopus»). Рвати адресу й так дозволяє xurl, тож проміжок лише шкодить.
\setlength{\biburlbigskip}{0mu}
% Без додаткового інтервалу між описами: розділ має йти суцільним текстом, як решта звіту. Так само в шаблоні бакалаврської.
\setlength{\bibitemsep}{0pt}
\setlength{\bibparsep}{0pt}

% ---------------------------- Лістинги програмного коду ---------------------------- Використано minted (через Pygments), а не listings: listings не вміє коректно обробляти багатобайтові символи й «висмикує» кириличні коментарі з коду, зліплюючи їх в один рядок. Потребує прапорця -shell-escape.
\usepackage{xcolor}
\usepackage[newfloat]{minted}
% labelsep=period дає «Рис. 2.1. Назва» замість «Рис. 2: Назва» — формат НУЛП.
\usepackage[labelsep=period]{caption}

% baselinestretch=1 обов'язковий: без нього лістинг успадковує півторашку основного тексту, і довгий фрагмент коду перестає вміщатися у сторінку. Усередині [H]-флоата це не помилка збірки, а нескінченний цикл порожніх сторінок — lab1 так намолотив 161 тисячу сторінок.
\setminted{
  baselinestretch=1,
  fontsize=\footnotesize,
  linenos=true,
  frame=single,
  framesep=4pt,
  breaklines=true,
  breakanywhere=true,
  tabsize=4,
  numbersep=7pt,
  bgcolor=gray!8,
}
\setminted[python]{python3=true}

\SetupFloatingEnvironment{listing}{name=Лістинг,listname=Перелік лістингів}

% ---------------------------- Нумерація та оформлення заголовків ----------------------------
\usepackage{titlesec}

\setcounter{secnumdepth}{3}
\renewcommand{\thesection}{\arabic{section}}
\renewcommand{\thesubsection}{\arabic{section}.\arabic{subsection}}
\renewcommand{\thesubsubsection}{\arabic{section}.\arabic{subsection}.\arabic{subsubsection}}

% Розділ по центру, 16 пт, великими літерами; слово «РОЗДІЛ» опущене, бо звіт не дисертація. Номер лишається: \numberwithin прив'язує до нього нумерацію рисунків («Рис. 2.1»), і без нумерованих розділів вона б зламалася. \MakeUppercase четвертим аргументом отримує текст заголовка, тому в джерелі заголовки пишуться звичайним регістром.
\titleformat{\section}[block]
  {\centering\bfseries\fontsize{16}{19.2}\selectfont}
  {\thesection.}{0.3em}{\MakeUppercase}
% Підрозділи — зліва з абзацного відступу, звичайним регістром.
\titleformat{\subsection}[block]{\bfseries}{\thesubsection.}{1em}{}
\titleformat{\subsubsection}[block]{\bfseries\itshape}{\thesubsubsection.}{1em}{}

% Незіркові \titlespacing — обов'язкові. Без них titlesec застосовує зірковий варіант, який прибирає абзацний відступ у першого абзацу після заголовка й перекриває indentfirst. Значення — типові для report.
\titlespacing{\section}{0pt}{20pt}{6pt}
\titlespacing{\subsection}{1.27cm}{6pt}{0pt}
\titlespacing{\subsubsection}{1.27cm}{6pt}{0pt}

\setlength{\parindent}{1.27cm}
\setlength{\parskip}{0pt}
\renewcommand{\arraystretch}{1.3}

% Переноси слів вимкнені — так у шаблоні НУЛП. Натомість абзац розганяється міжслівними пробілами, тому tolerance і emergencystretch підняті: без них рядки виходили б за праве поле.
\hyphenpenalty=10000
\exhyphenpenalty=10000
\tolerance=9999
\setlength{\emergencystretch}{3em}

% ---------------------------- Підписи до рисунків і таблиць ----------------------------
\captionsetup[figure]{name=Рис.,justification=centering,font={bf,it}}
\captionsetup[table]{name=Таблиця,justification=raggedleft,singlelinecheck=false,labelfont=bf,textfont=normalfont}

% Рисунки, таблиці й формули нумеруються в межах розділу: «Рис. 2.1», а не «Рис. 5». Так у шаблоні НУЛП, і так само виглядали приклади на консультації.
\numberwithin{equation}{section}
\numberwithin{figure}{section}
\numberwithin{table}{section}

% ---------------------------- Колонтитули ----------------------------
\usepackage{fancyhdr}
\pagestyle{fancy}
\fancyhf{}
\fancyfoot[R]{\thepage}
\renewcommand{\headrulewidth}{0pt}
\fancypagestyle{plain}{%
  \fancyhf{}
  \fancyfoot[R]{\thepage}
  \renewcommand{\headrulewidth}{0pt}
}

% titlepage.tex — титульний аркуш, спільний для всіх дисциплін.
% Потребує визначених у subject.tex команд:
%   \SubjectName, \LecturerName, \StudentName, \StudentGroup, \AcademicYear
%
% Використання у тексті: \MakeLabTitle{1}{Тема роботи}
%
% Рід роботи. Дисципліни з практичними роботами перевизначають його після
% % preamble.tex — спільна преамбула для всіх лабораторних робіт. Підключається одразу після \documentclass у кожному lab*/main.tex.
%
% ЗБІРКА: LuaLaTeX (не pdfLaTeX). Причина: пакет listings під pdfLaTeX не читає UTF-8 і зупиняє збірку на кириличних коментарях у лістингах програмного коду. LuaLaTeX обробляє UTF-8 нативно.
%
% ШРИФТ: Liberation Serif — метрично сумісний з Times New Roman і містить повний набір кириличних гліфів (TeX Gyre Termes кирилиці не має, через що український текст не відображався взагалі).

\usepackage{fontspec}
\setmainfont{Liberation Serif}
\setmonofont{DejaVu Sans Mono}[Scale=0.82]

\usepackage{polyglossia}
\setmainlanguage{ukrainian}
\setotherlanguage{english}

\usepackage{microtype}

% Поля, інтервал, відступи й заголовки нижче взяті з шаблону НУЛП, за яким написана бакалаврська робота (../bachelor-thesis/thesis/main.tex). Той самий формат приймає В.М. Шиманський і для магістерських, тож усі звіти тримаємо в ньому, а не в тому, що переказували на консультації.
\usepackage[left=25mm,right=15mm,top=15mm,bottom=23mm,footskip=1.5mm]{geometry}
\usepackage{setspace}
% Саме \setstretch{1.5}, а не \onehalfspacing: останній дає близько 1.25.
\setstretch{1.5}

\usepackage{graphicx}
\usepackage{float}
\usepackage{booktabs}
\usepackage{array}
\usepackage{tabularx}
\usepackage{longtable}
\usepackage{enumitem}
% labelindent вирівнює маркер списку з абзацним відступом — так у шаблоні НУЛП.
\setlist{nosep}
\setlist[itemize]{label=\textbullet,leftmargin=*,labelindent=1.27cm}
\setlist[enumerate]{leftmargin=*,labelindent=1.27cm}
\usepackage{indentfirst}
\usepackage{amsmath}

\usepackage{hyperref}
\hypersetup{colorlinks=true,linkcolor=black,urlcolor=blue,citecolor=black}

% xurl дозволяє \url розривати адресу в будь-якому місці, а не лише після скісної риски. Без нього довгий сегмент шляху (напр. «scopus---podrobnaya- instruktsiya-dlya-uchenogo») виходить за межі сторінки, бо дефіс для hyperref не є точкою розриву. Вантажиться після hyperref — перевизначає його \UrlBreaks.
\usepackage{xurl}

% ---------------------------- Список використаної літератури ----------------------------
% biblatex + biber зі стилем IEEE: посилання «[1]» у тексті, записи — за порядком першої згадки. Обрано IEEE, а не gost-numeric (ДСТУ 8302:2015), — так вирішив автор звітів після порівняння обох варіантів.
% Кожен звіт із джерелами сам підключає свій references.bib через \addbibresource і друкує список \printbibliography[heading=none] під власним \section. Звіт без джерел від завантаженого biblatex нічого не втрачає.
% Оточення списку — стандартне для стилю: мітка від лівого поля, текст опису вирівняний після неї. Не підганяємо його під абзацний відступ нумерованих переліків тексту (labelindent у \setlist вище) — список літератури є окремим розділом зі своїм стандартним оформленням.
\usepackage[backend=biber,style=ieee]{biblatex}
% biblatex за замовчуванням додає розтяжний проміжок після «/ . : ? =» в адресах, і з вимкненими переносами рядок вирівнюється розрідженням самої адреси («https : / / www . scopus»). Рвати адресу й так дозволяє xurl, тож проміжок лише шкодить.
\setlength{\biburlbigskip}{0mu}
% Без додаткового інтервалу між описами: розділ має йти суцільним текстом, як решта звіту. Так само в шаблоні бакалаврської.
\setlength{\bibitemsep}{0pt}
\setlength{\bibparsep}{0pt}

% ---------------------------- Лістинги програмного коду ---------------------------- Використано minted (через Pygments), а не listings: listings не вміє коректно обробляти багатобайтові символи й «висмикує» кириличні коментарі з коду, зліплюючи їх в один рядок. Потребує прапорця -shell-escape.
\usepackage{xcolor}
\usepackage[newfloat]{minted}
% labelsep=period дає «Рис. 2.1. Назва» замість «Рис. 2: Назва» — формат НУЛП.
\usepackage[labelsep=period]{caption}

% baselinestretch=1 обов'язковий: без нього лістинг успадковує півторашку основного тексту, і довгий фрагмент коду перестає вміщатися у сторінку. Усередині [H]-флоата це не помилка збірки, а нескінченний цикл порожніх сторінок — lab1 так намолотив 161 тисячу сторінок.
\setminted{
  baselinestretch=1,
  fontsize=\footnotesize,
  linenos=true,
  frame=single,
  framesep=4pt,
  breaklines=true,
  breakanywhere=true,
  tabsize=4,
  numbersep=7pt,
  bgcolor=gray!8,
}
\setminted[python]{python3=true}

\SetupFloatingEnvironment{listing}{name=Лістинг,listname=Перелік лістингів}

% ---------------------------- Нумерація та оформлення заголовків ----------------------------
\usepackage{titlesec}

\setcounter{secnumdepth}{3}
\renewcommand{\thesection}{\arabic{section}}
\renewcommand{\thesubsection}{\arabic{section}.\arabic{subsection}}
\renewcommand{\thesubsubsection}{\arabic{section}.\arabic{subsection}.\arabic{subsubsection}}

% Розділ по центру, 16 пт, великими літерами; слово «РОЗДІЛ» опущене, бо звіт не дисертація. Номер лишається: \numberwithin прив'язує до нього нумерацію рисунків («Рис. 2.1»), і без нумерованих розділів вона б зламалася. \MakeUppercase четвертим аргументом отримує текст заголовка, тому в джерелі заголовки пишуться звичайним регістром.
\titleformat{\section}[block]
  {\centering\bfseries\fontsize{16}{19.2}\selectfont}
  {\thesection.}{0.3em}{\MakeUppercase}
% Підрозділи — зліва з абзацного відступу, звичайним регістром.
\titleformat{\subsection}[block]{\bfseries}{\thesubsection.}{1em}{}
\titleformat{\subsubsection}[block]{\bfseries\itshape}{\thesubsubsection.}{1em}{}

% Незіркові \titlespacing — обов'язкові. Без них titlesec застосовує зірковий варіант, який прибирає абзацний відступ у першого абзацу після заголовка й перекриває indentfirst. Значення — типові для report.
\titlespacing{\section}{0pt}{20pt}{6pt}
\titlespacing{\subsection}{1.27cm}{6pt}{0pt}
\titlespacing{\subsubsection}{1.27cm}{6pt}{0pt}

\setlength{\parindent}{1.27cm}
\setlength{\parskip}{0pt}
\renewcommand{\arraystretch}{1.3}

% Переноси слів вимкнені — так у шаблоні НУЛП. Натомість абзац розганяється міжслівними пробілами, тому tolerance і emergencystretch підняті: без них рядки виходили б за праве поле.
\hyphenpenalty=10000
\exhyphenpenalty=10000
\tolerance=9999
\setlength{\emergencystretch}{3em}

% ---------------------------- Підписи до рисунків і таблиць ----------------------------
\captionsetup[figure]{name=Рис.,justification=centering,font={bf,it}}
\captionsetup[table]{name=Таблиця,justification=raggedleft,singlelinecheck=false,labelfont=bf,textfont=normalfont}

% Рисунки, таблиці й формули нумеруються в межах розділу: «Рис. 2.1», а не «Рис. 5». Так у шаблоні НУЛП, і так само виглядали приклади на консультації.
\numberwithin{equation}{section}
\numberwithin{figure}{section}
\numberwithin{table}{section}

% ---------------------------- Колонтитули ----------------------------
\usepackage{fancyhdr}
\pagestyle{fancy}
\fancyhf{}
\fancyfoot[R]{\thepage}
\renewcommand{\headrulewidth}{0pt}
\fancypagestyle{plain}{%
  \fancyhf{}
  \fancyfoot[R]{\thepage}
  \renewcommand{\headrulewidth}{0pt}
}

\input{../../../../titlepage}
: \renewcommand{\WorkKind}{ПРАКТИЧНА РОБОТА}.
\newcommand{\WorkKind}{ЛАБОРАТОРНА РОБОТА}

\newcommand{\MakeLabTitle}[2]{%
\begin{titlepage}
% Титулка — не суцільний текст, а фіксована верстка з явних \\[...] проміжків,
% тож міжрядковий інтервал основного тексту (\setstretch{1.5} у preamble.tex)
% до неї не застосовний: він розтягував кожен із ~15 рядків на 20 %, і блок
% «Виконав» разом із «Львів --- рік» переповнював сторінку. Інтервал локальний
% для оточення titlepage, тому на решту документа це не впливає.
\setstretch{1}
\begin{center}
  {\large\bfseries МІНІСТЕРСТВО ОСВІТИ І НАУКИ УКРАЇНИ}\\[6pt]
  {\large\bfseries Національний університет «Львівська політехніка»}\\[6pt]
  {\large\bfseries Інститут комп'ютерних наук та інформаційних технологій}\\[6pt]
  {\large\bfseries Кафедра «Системи штучного інтелекту»}\\[1.2cm]

  \includegraphics[width=5cm]{../../../../logo}\\[1.2cm]

  {\Large\bfseries \WorkKind{} №#1}\\[0.5cm]
  з дисципліни\\[4pt]
  «\SubjectName»\\[0.6cm]
  на тему\\[4pt]
  \textit{#2}
\end{center}

% Відступ до блоку «Виконав» гнучкий, а не фіксований (було 2 cm). Фіксований
% разом із довгим родом роботи чи темою на два рядки переповнював сторінку, і
% «Львів --- рік» зі свого \vfill з'їжджав на окрему другу сторінку. \vfill
% стискається до нуля, коли місця бракує, тож переповнення стало неможливим.
\vfill

\begin{flushright}
  {\bfseries Виконав:}\\
  студент групи \StudentGroup\\
  \StudentName\\[1cm]
  {\bfseries Викладач:}\\
  \LecturerName\\
\end{flushright}

\vfill
\begin{center}
  Львів --- \AcademicYear
\end{center}
\end{titlepage}
\clearpage
}


% Метадані дисципліни. Використовуються титульним аркушем (titlepage.tex).

\newcommand{\SubjectName}{Науковий процес та професійна науково-дослідна робота}
% Два рядки: посада з науковим ступенем, далі прізвище. `\\` розриває рядок усередині flushright-блоку титульного аркуша. Посаду (професор, не доцент) підтвердив сам викладач.
\newcommand{\LecturerName}{проф. каф. СШІ д.т.н.,\\ Ізонін І.В.}
\newcommand{\StudentName}{Галич Максим}
\newcommand{\StudentGroup}{КНСШ-13}
\newcommand{\AcademicYear}{2026}

% Same work kind as in practice1: the document reports on a practice, it is not the assignment.
\renewcommand{\WorkKind}{ЗВІТ\\[2pt]\normalfont\large про виконання практичної роботи}
% Screenshots are hand-made input that cannot be regenerated, so they live next to the text and are committed.
\graphicspath{{assets/}}
\addbibresource{references.bib}

% Full-width screenshot figure, styled as in practice2 and practice3.
\newcommand{\Screenshot}[4][]{%
  \begin{figure}[H]
    \centering
    \includegraphics[width=\linewidth,#1]{#2}
    \caption{#3}
    \label{#4}
  \end{figure}
}
% Half-width screenshot with a letter under it, two per row inside one figure. The preamble has no subcaption, and the figure's own caption names each letter.
\newcommand{\Tile}[2]{%
  \begin{minipage}[t]{0.49\linewidth}
    \centering
    \includegraphics[width=\linewidth]{#1}\\[2pt]
    {\small #2}
  \end{minipage}%
}

\begin{document}

\MakeLabTitle{4}{Вибір надійного наукового видання для опублікування результатів наукових досліджень}

\textbf{Мета роботи.} Навчитися шукати, обирати та обґрунтовувати вибір періодичного наукового видання для опублікування результатів наукових досліджень (українські та закордонні наукові видання).

\section{Постановка завдання}

Здійснити пошук, вибір та обґрунтування вибору наукового періодичного видання (українського та закордонного) для опублікування результатів наукових досліджень~\cite{guide}. Зокрема:

\begin{enumerate}
\item написати назву, анотацію та ключові слова статті англійською та українською мовами;

\item за допомогою сервісу Journal Finder видавництва Elsevier~\cite{journalfinder} знайти журнали, релевантні статті англійською мовою, навести щонайменше п'ять журналів з їхніми основними показниками (імпакт-фактор, квартиль, тривалість рецензування, відсоток прийнятих статей тощо), обрати один із них і обґрунтувати вибір;

\item за допомогою Реєстру наукових фахових видань України~\cite{nfv} знайти журнали з технічних наук (спеціальність «Комп'ютерні науки»), релевантні статті українською мовою, навести щонайменше п'ять журналів з їхніми основними показниками (категорія видання, індексація в наукометричних базах, вартість публікації, кількість випусків на рік тощо), обрати один із них і обґрунтувати вибір;

\item описати проведену роботу, проілюструвати її рисунками й таблицями з посиланнями на них у тексті;

\item зробити висновок: що зроблено, як саме і як це допоможе під час написання магістерської роботи.
\end{enumerate}

\section{Опис проведеної роботи}

\subsection{Стаття для публікації}

Опублікованих статей автор не має: на момент виконання практичної роботи №1 база Scopus не містила жодної його праці. Тому, як дозволяють методичні рекомендації~\cite{guide}, використано нову статтю за темою магістерської роботи «Удосконалення методу синтезу мовлення з керуванням голосом за текстовим описом для української мови». Стаття запланована, тож анотація описує проблему, запропонований метод і спосіб його оцінювання, але не наводить результатів.

Назва, анотація та ключові слова статті англійською мовою:

\begin{english}
\begin{description}
\item[Title:] Text-Described Voice Control for Ukrainian Speech Synthesis

\item[Abstract:] Speech synthesis models with voice control by text description let a user set the voice with a natural-language prompt, such as ``a calm low female voice speaking slowly'', instead of a reference recording. Current models of this type are trained mostly on English and Chinese data. For Ukrainian, a low-resource language, there is no corpus of speech paired with voice descriptions and no model that follows such descriptions. This paper proposes a method of text-described voice control for Ukrainian speech synthesis. The method builds a corpus of Ukrainian speech with automatically generated descriptions of speaker attributes (gender, age, pitch, speaking rate, emotion) and adapts a multilingual prompt-conditioned speech synthesis model to this corpus. The evaluation measures how well the synthesized speech matches the description and how intelligible and natural it is, with objective metrics and a listening test. The corpus construction procedure can be reused for other low-resource languages.

\item[Keywords:] speech synthesis; text-to-speech; voice control; text description; Ukrainian language; low-resource language.
\end{description}
\end{english}

Назва, анотація та ключові слова статті українською мовою:

\begin{description}
\item[Назва:] Синтез українського мовлення з керуванням голосом за текстовим описом

\item[Анотація:] Моделі синтезу мовлення з керуванням голосом за текстовим описом дозволяють задати голос природномовною підказкою, наприклад «спокійний низький жіночий голос, повільний темп», замість еталонного запису. Наявні моделі цього типу навчено переважно на англійських і китайських даних. Для української мови, яка належить до мов з обмеженими ресурсами, немає ні корпусу мовлення з описами голосу, ні моделі, що виконує такі описи. У статті запропоновано метод синтезу українського мовлення з керуванням голосом за текстовим описом. Метод формує корпус українського мовлення з автоматично згенерованими описами характеристик мовця (стать, вік, висота тону, темп, емоція) та адаптує до нього багатомовну модель синтезу мовлення, керовану текстовою підказкою. Оцінювання вимірює відповідність синтезованого мовлення описові, його розбірливість і природність за об'єктивними метриками та під час аудитивного тесту. Процедуру формування корпусу можна застосувати й до інших мов з обмеженими ресурсами.

\item[Ключові слова:] синтез мовлення; керування голосом; текстовий опис голосу; українська мова; мова з обмеженими ресурсами.
\end{description}

\subsection{Пошук закордонного журналу в Journal Finder}

Пошук закордонного журналу виконано в сервісі Journal Finder видавництва Elsevier~\cite{journalfinder}. Сервіс має два режими: пошук за анотацією статті (\textit{Match my abstract}) і пошук за назвою журналу, предметною галуззю чи тематикою (\textit{Search by journal title, subject area, or aims \& scope}). Другий режим зіставляє запит із назвами й описами журналів, а не з текстом статті, тому за назвою статті він не знайшов жодного журналу (рисунок~\ref{fig:jf-title}).

\Screenshot{jf-title-search}{Пошук за назвою статті в режимі пошуку журналів: жодного результату}{fig:jf-title}

У режимі \textit{Match my abstract} у поле запиту вставлено анотацію англомовної статті. Сервіс знайшов 40 журналів видавництва Elsevier, упорядкованих за відповідністю статті (рисунок~\ref{fig:jf-abstract}). Для кожного журналу видача наводить імпакт-фактор, CiteScore, середній час від подання статті до першого рішення редакції та від прийняття статті до її публікації, а також умови публікації: вартість відкритого доступу або безкоштовну публікацію за передплатною моделлю.

\Screenshot{jf-abstract}{Видача Journal Finder за анотацією статті: 40 журналів, перший — Speech Communication}{fig:jf-abstract}

Сервіс упорядковує видачу за подібністю тексту анотації до статей журналу, а не за його тематикою, тому кожен журнал перевірено за предметними галузями (\textit{Subject Areas}) та описом тематики (\textit{Journal Scope}), які показує розгорнута картка журналу (\textit{More about this journal}). Переглянуто 12 журналів видачі (рисунок~\ref{fig:jf-list}), з них для порівняння відібрано п'ять, тематика яких охоплює обробку мовлення або методи штучного інтелекту для мовних даних.

\begin{figure}[H]
  \centering
  \Tile{jf-list-1}{а)}\hfill\Tile{jf-list-2}{б)}
  \par\medskip
  \Tile{jf-list-3}{в)}\hfill\Tile{jf-list-4}{г)}
  \caption{Журнали з видачі Journal Finder: а) Speech Communication, Data in Brief, Computer Speech and Language; б) Natural Language Processing Journal, Journal of Voice, Applied Acoustics; в) SoftwareX, Computers \& Electrical Engineering, Ampersand; г) Expert Systems with Applications, Heliyon, Applied Corpus Linguistics}
  \label{fig:jf-list}
\end{figure}

Решту журналів відхилено з таких причин:

\begin{itemize}
\item Journal of Voice — назва збігається з темою статті, проте це клінічний журнал: його предметні галузі — отоларингологія та лінгвістика, а сам журнал описує себе як видання з медицини голосу (\textit{voice medicine}), тобто про порушення голосу, а не про його синтез (рисунок~\ref{fig:jf-voice});
\item Data in Brief і SoftwareX публікують описи наборів даних і програмного забезпечення, а не нові методи; Data in Brief може підійти для окремої статті з описом корпусу, але не для статті про метод синтезу;
\item Heliyon — мультидисциплінарний журнал, Ampersand і Applied Corpus Linguistics — лінгвістичні, Computers \& Electrical Engineering — загальноінженерний; тематика жодного з них не охоплює синтезу мовлення.
\end{itemize}

\Screenshot{jf-journal-of-voice}{Розгорнута картка Journal of Voice: предметні галузі та тематика журналу}{fig:jf-voice}

Розгорнуті картки відібраних журналів наведено на рисунках~\ref{fig:jf-sc} та~\ref{fig:jf-cards}. Окрім показників видачі, картка містить ISSN журналу, а для двох журналів із дванадцяти — Speech Communication і Journal of Voice — ще й частку прийнятих статей (\textit{Acceptance rate}); для решти журналів цей показник Journal Finder не наводить.

\Screenshot{jf-speech-communication}{Розгорнута картка Speech Communication: ISSN, частка прийнятих статей, предметні галузі}{fig:jf-sc}

\begin{figure}[H]
  \centering
  \Tile{jf-csl}{а)}\hfill\Tile{jf-applied-acoustics}{б)}
  \par\medskip
  \Tile{jf-eswa}{в)}\hfill\Tile{jf-nlp-journal}{г)}
  \caption{Розгорнуті картки журналів: а) Computer Speech and Language; б) Applied Acoustics; в) Expert Systems with Applications; г) Natural Language Processing Journal}
  \label{fig:jf-cards}
\end{figure}

Квартиль журналу Journal Finder не показує, тому його перевірено за ISSN у SCImago Journal \& Country Rank~\cite{scimago}. Сервіс обчислює показник SJR за даними бази Scopus і поряд з ним наводить квартиль журналу (рисунок~\ref{fig:sjr}).

\begin{figure}[H]
  \centering
  \Tile{sjr-speech-communication}{а)}\hfill\Tile{sjr-csl}{б)}
  \par\medskip
  \Tile{sjr-applied-acoustics}{в)}\hfill\Tile{sjr-eswa}{г)}
  \par\medskip
  \Tile{sjr-nlp-journal}{д)}
  \caption{Показник SJR 2025 і квартиль журналів у SCImago: а) Speech Communication; б) Computer Speech and Language; в) Applied Acoustics; г) Expert Systems with Applications; д) Natural Language Processing Journal}
  \label{fig:sjr}
\end{figure}

Показники п'яти відібраних журналів зведено в таблиці~\ref{tab:journals}. Назви журналів у заголовку таблиці скорочено: SC — Speech Communication, CSL — Computer Speech and Language, AA — Applied Acoustics, ESWA — Expert Systems with Applications, NLPJ — Natural Language Processing Journal.

\begin{table}[H]
\caption{Показники журналів, відібраних у Journal Finder}
\label{tab:journals}
\small
\begin{tabularx}{\linewidth}{|X|c|c|c|c|c|}
\hline
\textbf{Показник} & \textbf{SC} & \textbf{CSL} & \textbf{AA} & \textbf{ESWA} & \textbf{NLPJ} \\
\hline
ISSN & 0167-6393 & 0885-2308 & 0003-682X & 0957-4174 & 2949-7191 \\
\hline
Імпакт-фактор & 2,9 & 3,0 & 4,3 & 9,4 & — \\
\hline
CiteScore & 7,4 & 10,9 & 8,3 & 17,0 & 9,1 \\
\hline
SJR 2025 & 0,600 & 0,702 & 0,855 & 1,939 & 1,266 \\
\hline
Квартиль (SCImago) & Q1 & Q2 & Q1 & Q1 & Q1 \\
\hline
H-індекс & 123 & 94 & 122 & 315 & 16 \\
\hline
До першого рішення, днів & 49 & 3 & 5 & 5 & 10 \\
\hline
Від прийняття до публікації, днів & 7 & 12 & 9 & 7 & 5 \\
\hline
Частка прийнятих статей & 19\,\% & — & — & — & — \\
\hline
Модель публікації & гібридна & гібридна & гібридна & гібридна & лише OA \\
\hline
Вартість публікації у відкритому доступі, \$ & 3\,280 & 3\,240 & 4\,070 & 3\,630 & 915\textsuperscript{*} \\
\hline
\end{tabularx}

\medskip
\footnotesize\textsuperscript{*} Зі знижкою для статей, поданих до 31.12.2026; стандартна вартість — 1\,830~\$.
\end{table}

Для публікації статті обрано журнал \textbf{Speech Communication}. Вибір обґрунтовано так:

\begin{enumerate}
\item Журнал найкраще відповідає тематиці статті: це профільне видання з обробки мовлення, яке видають Європейська асоціація обробки сигналів (EURASIP) та Міжнародна асоціація мовленнєвої комунікації (ISCA), а серед його предметних галузей є обробка сигналів і лінгвістика (рисунок~\ref{fig:jf-sc}). Сервіс Journal Finder поставив його першим у видачі за анотацією статті.
\item Журнал належить до першого квартиля (Q1) за даними SCImago, тобто входить до чверті найвпливовіших журналів своєї предметної категорії за показником SJR.
\item Частка прийнятих статей становить 19\,\%: журнал приймає менше ніж кожну п'яту подану статтю, що свідчить про суворе рецензування. Серед п'яти відібраних журналів це єдиний, для якого цей показник оприлюднено.
\item Журнал має гібридну модель: статтю можна опублікувати безкоштовно за передплатною моделлю, а відкритий доступ (3\,280~\$) — сплатити за наявності фінансування.
\item Від прийняття статті до її публікації минає в середньому 7 днів. Найдовший серед кандидатів час до першого рішення (49 днів) — недолік журналу, проте для статті без жорсткого строку публікації він прийнятний.
\end{enumerate}

Інші кандидати поступаються обраному журналу:

\begin{itemize}
\item Computer Speech and Language тематично майже рівноцінний — у його тематиці прямо згадано синтез мовлення, а перше рішення редакції він ухвалює за 3 дні, — проте належить до другого квартиля. Його обрано запасним варіантом, якщо Speech Communication відхилить статтю;
\item Expert Systems with Applications має найвищі показники (імпакт-фактор 9,4, SJR 1,939), але дуже широку тематику — прикладні інтелектуальні системи для промисловості, фінансів і менеджменту. Стаття про синтез мовлення в ньому ризикує бути відхиленою як така, що не відповідає тематиці, і не дійде до фахівців з обробки мовлення;
\item Applied Acoustics належить до першого квартиля, але присвячений інженерній акустиці (шумозаглушення, будівельна та підводна акустика) і вимагає експериментальної перевірки на практичній акустичній задачі;
\item Natural Language Processing Journal найдешевший, проте публікує лише у відкритому доступі, тобто безкоштовної публікації не передбачає. Журнал індексується в Scopus лише з 2024~р., ще не має імпакт-фактора (H-індекс 16), а його тематика зосереджена на обробці тексту.
\end{itemize}

\subsection{Пошук українського фахового видання}

Методичні рекомендації пропонують шукати українське фахове видання в Реєстрі наукових фахових видань України на сайті УкрІНТЕІ~\cite{guide,nfv}. На момент виконання роботи цей реєстр недоступний: його адреса перенаправляє на головну сторінку УкрІНТЕІ (рисунок~\ref{fig:nfv}). Перелік наукових фахових видань України тепер наводить інтерактивний дашборд «Фахові видання України» на порталі даних Міністерства освіти і науки~\cite{mondash}.

\Screenshot{nfv-redirect}{Головна сторінка УкрІНТЕІ, на яку перенаправляє адреса реєстру з методичних рекомендацій}{fig:nfv}

Дашборд містить 1\,302 видання, і всі вони належать до категорії Б (рисунок~\ref{fig:mon}). Видання категорії А, тобто ті, що індексуються в базах Web of Science Core Collection або Scopus~\cite{guide}, дашборд обіцяє опублікувати згодом, тому українські журнали з індексацією в цих базах знайти в ньому поки неможливо.

\Screenshot{mon-dashboard}{Дашборд «Фахові видання України»: 1\,302 видання категорії Б, категорію А ще не опубліковано}{fig:mon}

Спеціальності в дашборді наведено за новим переліком, у якому спеціальності «Комп'ютерні науки» відповідає код F3. Фільтр за цією спеціальністю залишає 78 видань від 55 засновників (рисунок~\ref{fig:mon-f3}). З них 51 видання має діамантовий відкритий доступ, за якого ні автор не платить за публікацію, ні читач за доступ до статті, а 13 видань входять до каталогу журналів відкритого доступу DOAJ. Для кожного видання таблиця дашборда наводить засновника, спеціальності, періодичність і ISSN (рисунок~\ref{fig:mon-f3-table}).

\Screenshot{mon-f3-summary}{Видання за спеціальністю F3 «Комп'ютерні науки»: 78 видань категорії Б}{fig:mon-f3}

\Screenshot{mon-f3-table}{Таблиця видань за спеціальністю F3 у дашборді}{fig:mon-f3-table}

З 78 видань за назвою та тематикою відібрано п'ять, найближчих до теми статті. Для кожного з них категорію, періодичність, наявність діамантового доступу та входження до DOAJ перевірено пошуком за ISSN у дашборді (рисунок~\ref{fig:mon-five}). Індексацію в наукометричних базах і вартість публікації дашборд не наводить, тому їх взято із сайтів журналів~\cite{jai,ujit,sisn,nltu,itce}.

\begin{figure}[H]
  \centering
  \Tile{mon-jai}{а)}\hfill\Tile{mon-ujit}{б)}
  \par\medskip
  \Tile{mon-isn}{в)}\hfill\Tile{mon-nltu}{г)}
  \par\medskip
  \Tile{mon-itce}{д)}
  \caption{Пошук відібраних видань за ISSN у дашборді: а) «Штучний інтелект»; б) «Український журнал інформаційних технологій»; в) «Вісник НУ «Львівська політехніка». Інформаційні системи та мережі»; г) «Науковий вісник НЛТУ України»; д) «Інформаційні технології та комп'ютерна інженерія»}
  \label{fig:mon-five}
\end{figure}

Показники відібраних видань зведено в таблиці~\ref{tab:ukr}. Усі п'ять належать до категорії Б і кластера «Інформаційні технології та електроніка», тому ці показники в таблиці не повторюються.

\begin{table}[H]
\caption{Показники українських фахових видань за спеціальністю F3}
\label{tab:ukr}
\small
\begin{tabularx}{\linewidth}{|>{\raggedright\arraybackslash}X|c|c|c|>{\raggedright\arraybackslash}X|}
\hline
\textbf{Видання, засновник, ISSN} & \textbf{Вип./рік} & \textbf{Діам. доступ} & \textbf{DOAJ} & \textbf{Індексація} \\
\hline
«Штучний інтелект», ДУІКТ, 2710-1673 & 4 & ні & ні & Google Scholar, Index Copernicus, Ulrich's \\
\hline
«Український журнал інформаційних технологій», НУ «Львівська політехніка», 2707-1898 & 2 & так & ні & Index Copernicus, Crossref, Google Scholar, НБУВ \\
\hline
«Вісник НУ «Львівська політехніка». Інформаційні системи та мережі», 2524-065X & 2 & так & ні & Google Scholar, Index Copernicus \\
\hline
«Науковий вісник НЛТУ України», НЛТУ, 1994-7836 & 6 & так & ні & Index Copernicus, Crossref, Scilit, WorldCat, НБУВ \\
\hline
«Інформаційні технології та комп'ютерна інженерія», ВНТУ, 1999-9941 & 3 & так & так & DOAJ, Google Scholar, Ulrichsweb, OUCI, CORE, J-Gate, WorldCat, НБУВ \\
\hline
\end{tabularx}
\end{table}

Видання з діамантовим доступом публікують статті безкоштовно; журнал «Інформаційні технології та комп'ютерна інженерія» прямо зазначає це на своєму сайті~\cite{itce}. Журнал «Штучний інтелект» діамантового доступу не має, а вартості публікації на його сайті не оприлюднено; за даними дашборда, від оплати звільнено лише аспірантів. Жодне з п'яти видань не індексується в Scopus чи Web of Science, що узгоджується з їхньою категорією Б.

Для публікації статті українською мовою обрано \textbf{«Науковий вісник НЛТУ України»}. Вибір обґрунтовано так:

\begin{enumerate}
\item Журнал публікує праці, суміжні з темою статті: 2026~р. у ньому вийшла стаття Біласа і Різника~\cite{bilas2026} про генерування штучної мови на основі фонетичних даних, результати якої оцінено аудитивним тестом за участі 32 респондентів. Отже, редакція приймає праці з фонетичних даних, оцінені тим самим способом — аудитивним тестом, — що й запланований у статті.
\item Журнал має діамантовий відкритий доступ, тобто публікація безкоштовна, а стаття вільно доступна читачам.
\item Журнал виходить 6 разів на рік — найчастіше серед відібраних видань, тож від прийняття статті до її виходу минає найменше часу.
\item Статтям присвоюється DOI через Crossref, а журнал індексують Index Copernicus, Scilit і WorldCat~\cite{nltu} — ширше, ніж три з чотирьох інших видань.
\end{enumerate}

Інші видання поступаються обраному:

\begin{itemize}
\item «Інформаційні технології та комп'ютерна інженерія» має найширшу індексацію і єдиний з п'яти входить до DOAJ, проте присвячений інформаційним технологіям і комп'ютерній інженерії загалом і виходить 3 рази на рік. Його обрано запасним варіантом;
\item «Штучний інтелект» тематично найближчий — серед його напрямів є природномовні системи та системи розпізнавання образів~\cite{jai}, — проте публікація в ньому платна, а її вартість не оприлюднено;
\item «Вісник НУ «Львівська політехніка». Інформаційні системи та мережі» охоплює комп'ютерну лінгвістику~\cite{sisn}, але виходить лише двічі на рік, статті приймає до 1~березня та 10~вересня, а індексується лише в Google Scholar та Index Copernicus;
\item «Український журнал інформаційних технологій» має загальну тематику інформаційних технологій і виходить двічі на рік.
\end{itemize}

\section{Результати}

У результаті роботи для статті за темою магістерської роботи обрано два наукові періодичні видання (таблиця~\ref{tab:choice}): закордонне — для англомовної статті й українське фахове — для україномовної. Для кожного визначено й запасний варіант на випадок відхилення статті.

\begin{table}[H]
\caption{Обрані видання для публікації статті}
\label{tab:choice}
\small
\begin{tabularx}{\linewidth}{|l|X|X|}
\hline
\textbf{Ознака} & \textbf{Закордонне видання} & \textbf{Українське видання} \\
\hline
Обране видання & Speech Communication & Науковий вісник НЛТУ України \\
\hline
Запасне видання & Computer Speech and Language & Інформаційні технології та комп'ютерна інженерія \\
\hline
Засіб пошуку & Elsevier Journal Finder, SCImago & дашборд МОН «Фахові видання України» \\
\hline
Категорія або квартиль & Q1 (SCImago) & категорія Б \\
\hline
Вартість публікації & безкоштовно за передплатною моделлю, 3\,280~\$ за відкритий доступ & безкоштовно (діамантовий доступ) \\
\hline
Швидкість & 49 днів до першого рішення, 7 днів від прийняття до публікації & 6 випусків на рік \\
\hline
\end{tabularx}
\end{table}

\section{Висновки}

У ході практичної роботи сформульовано назву, анотацію та ключові слова статті за темою магістерської роботи англійською та українською мовами, за анотацією знайдено закордонні журнали в сервісі Journal Finder, їхні квартилі перевірено в SCImago, а українські фахові видання відібрано в дашборді МОН. Для кожної мови обрано журнал і обґрунтовано вибір.

\begin{enumerate}[ref=\arabic*]
\item Journal Finder швидко звужує вибір із тисяч журналів до кількох десятків, проте впорядковує їх за подібністю тексту, а не за тематикою. Тому кожен журнал потрібно перевіряти за предметними галузями й описом тематики: так, Journal of Voice, попри назву, виявився клінічним журналом. Крім того, сервіс охоплює лише журнали Elsevier, тож для повного огляду варто скористатися й аналогічними сервісами інших видавців.
\item Жоден показник окремо не визначає найкращого журналу. Computer Speech and Language має вищий CiteScore, ніж Speech Communication, але нижчий квартиль; Expert Systems with Applications має найвищі показники серед кандидатів, але його тематика не відповідає статті. Вибір має поєднувати відповідність тематиці, квартиль, модель публікації та строки розгляду.
\item Частку прийнятих статей, яку методичні рекомендації називають серед основних показників, оприлюднюють рідко: Journal Finder навів її лише для двох журналів із дванадцяти.
\item Реєстр наукових фахових видань, на який посилаються методичні рекомендації, більше недоступний; його замінив дашборд МОН, у якому ще немає видань категорії А. Серед видань категорії Б за спеціальністю «Комп'ютерні науки» безкоштовна публікація — звичайна умова (51 із 78), тож між ними вибір визначають тематика, індексація та частота виходу.
\item Для магістерської роботи результат роботи — готовий план публікацій: основний результат — англомовною статтею до Speech Communication, україномовну статтю — до «Наукового вісника НЛТУ України», із запасними виданнями для кожної. Строки розгляду обраних журналів, зокрема 49 днів до першого рішення в Speech Communication, слід заздалегідь закласти в графік роботи над магістерською.
\end{enumerate}

\section{Список використаної літератури}

\printbibliography[heading=none]

\end{document}
